% Lecture 05: six pending source questions, synchronized with docs/lectures/05.md.
% Build from the repository root: make questions LECTURE=05
% Store evidence crops in the new project's images/ folder.
\documentclass[a4paper,12pt]{article}
\usepackage[margin=19mm]{geometry}
\usepackage{amsmath,amssymb,graphicx,xcolor}
\usepackage{fontspec}
\usepackage{xepersian}
\settextfont{XB Roya}
\setlatintextfont{Latin Modern Roman}
\definecolor{ink}{HTML}{173D75}
\graphicspath{{images/}}
\setlength{\parindent}{0pt}
\setlength{\parskip}{5pt}
\linespread{1.12}
\newcommand{\qtitle}[2]{{\Large\bfseries\color{ink} پرسش \lr{#1}: #2}\par}
\newcommand{\answer}{\par\medskip\textbf{پاسخ یا صورت مورد تأیید و تاریخ:}\par\vspace{18mm}\hrule\medskip}

\newcommand{\sourcepage}[1]{%
  \textbf{نشانی منبع:} نسخهٔ \lr{primary}، صفحهٔ #1.\par
  {\footnotesize\noindent\lr{\texttt{sources/lectures/05 Maximal submodule.pdf}}\par}}
\newcommand{\pending}{\textbf{وضعیت:} در انتظار پاسخ؛ تا تاریخ \lr{2026-10-25}
  پاسخی ثبت نشده است.\par}
\newcommand{\evidence}[1]{\begin{center}
  \includegraphics[width=\linewidth,height=72mm,keepaspectratio]{#1}
  \end{center}}

\begin{document}
{\LARGE\bfseries پرسش‌های ابهام جلسهٔ ۰۵}\par
{\large زیرمدول‌های ماکسیمال؛ جزوهٔ جبر پیشرفتهٔ دکتر خسروی}\par
این پرونده شش پرسش حل‌نشده را گرد می‌آورد. شماره‌ها، شمارهٔ یک‌مبنای صفحه در
فایل پی‌دی‌اف هستند، نه شمارهٔ دست‌نویس. پیشنهادها هنوز صورت قطعی اسلایدها نیستند.
برای هر پرسش، صورت مورد تأیید را همراه با شناسهٔ آن مشخص کنید.
\medskip\hrule\medskip

\qtitle{L05-Q01}{قید سره‌بودن ایده‌آل}
\pending
\sourcepage{۲، نتیجهٔ بالای صفحه، پس از پایان برهان قضیهٔ وجود زیرمدول ماکسیمال}
\textbf{عبارت منبع:} در حلقهٔ جابجایی و یک‌دار، «هر ایده‌آل» در یک ایده‌آل
ماکسیمال مشمول است. سپس از $R=(1)$ و تولید متناهی استفاده می‌شود.
\evidence{L05-Q01-ideals.png}
\textbf{خوانش‌های ممکن:} عبارت «هر ایده‌آل» به‌صورت لفظی؛ یا
«هر ایده‌آل \textbf{سره}» به‌عنوان اصلاح پیشنهادی.

\textbf{دلیل پرسش:} قضیهٔ قبل دربارهٔ زیرمدول‌های سره است؛ صورت لفظی، خود $R$
را هم شامل می‌شود، در حالی که $R$ در ایده‌آلی ماکسیمال مشمول نیست.

\textbf{پرسش:} آیا صورت اسلاید چنین باشد: «در حلقهٔ جابجایی و یک‌دار، هر ایده‌آل
سره در یک ایده‌آل ماکسیمال مشمول است»؟
\answer
\newpage

\qtitle{L05-Q02}{علامت سؤال کنار قید طول متناهی}
\pending
\sourcepage{۲، تذکرِ میانهٔ پایین صفحه، پس از مثال میدان}
\textbf{عبارت منبع:} «در مورد مدول‌ها (با طول متناهی؟)، مدول‌های ساده بلوک‌های
سازندهٔ آنها هستند؛ چیزی شبیه به گروه‌های ساده در نظریهٔ گروه‌ها.»
علامت سؤال در خود دست‌نویس دیده می‌شود.
\evidence{L05-Q02-length.png}
\textbf{خوانش‌های ممکن:}
\begin{enumerate}
    \item قید «با طول متناهی» فرض مورد نظر است و در اسلاید بدون علامت سؤال نوشته شود.
    \item این عبارت یادداشت پرسشیِ نویسنده است و همان قید و علامت سؤال، با توضیح
    اینکه از منبع است، حفظ شود.
\end{enumerate}
\textbf{دلیل پرسش:} تعبیرِ عوامل سادهٔ یک زنجیرهٔ ترکیب متناهی به قید طول متناهی
وابسته است؛ یادداشت پرسشی نباید به ادعایی بی‌قید تبدیل شود.

\textbf{پرسش:} کدام صورت را تأیید می‌کنید: بیان قید طول متناهی، یا حفظِ آشکارِ
علامت سؤالِ منبع؟
\answer
\newpage

\qtitle{L05-Q03}{فرض اول‌بودن عدد $p$}
\pending
\sourcepage{۲، پایین صفحه، دو مثالِ $\mathbb{Z}_p$}
\textbf{عبارت منبع:} در ادامهٔ مثال‌های مدول ساده، این دو مورد آمده‌اند:
\begin{itemize}
    \item $\mathbb{Z}_p$ به عنوان $\mathbb{Z}_p$-مدول؛
    \item $\mathbb{Z}_p$ به عنوان $\mathbb{Z}$-مدول.
\end{itemize}
در کنار این دو مثال، اول‌بودن $p$ تصریح نشده است.
\evidence{L05-Q03-prime.png}
\textbf{خوانش‌های ممکن:} $p$ از ابتدا نماد عددی اول است؛ یا شرط اول‌بودن در
این دو مثال افتاده و باید به‌صراحت افزوده شود.

\textbf{دلیل پرسش:} ساده‌بودن این دو مدول به اول‌بودن $p$ وابسته است؛
برای پیمانه‌ای مرکب، زیرمدول غیربدیهی وجود دارد.

\textbf{پرسش:} آیا برای هر دو مثال، عبارت «فرض کنید $p$ عددی اول باشد» را صریحاً
در اسلاید بیاورم؟
\answer
\newpage

\qtitle{L05-Q04}{شاخص مجموع مستقیم در نتیجه}
\pending
\sourcepage{۸، نتیجهٔ بالای صفحه، پس از برهان طولانی تجزیه}
\textbf{عبارت منبع:} از $M=\sum_{i\in I}M_i$ با $M_i$های ساده، وجود
$J\subseteq I$ بیان شده، اما فرمول نتیجه چنین نوشته شده است:
\[
    M=\bigoplus_{i\in I}M_i.
\]
\evidence{L05-Q04-indices.png}
\textbf{خوانش‌های ممکن:} شاخص $I$ به‌صورت لفظی؛ یا شاخص $J$ که با زیرخانوادهٔ
انتخاب‌شده در قضیهٔ قبل سازگار است.

\textbf{دلیل پرسش:} لازم نیست تمام خانوادهٔ اولیه مستقل باشد. برهان منبع فقط
می‌گوید در قضیهٔ قبل $K=(0)$ قرار دهید؛ آن قضیه یک زیرخانواده با شاخص $J$ می‌دهد.

\textbf{پرسش:} آیا فرمول نتیجه با اصلاح شاخص، چنین نوشته شود؟
\[
    M=\bigoplus_{j\in J}M_j.
\]
\answer
\newpage

\qtitle{L05-Q05}{قید سره‌بودن زیرمدول در مثال گویاها}
\pending
\sourcepage{۸، آخرین مثالِ پایین صفحه و سطرِ ناتمامِ پس از آن}
\textbf{عبارت منبع:} «در $\mathbb{Q}$ به عنوان $\mathbb{Z}$-مدول، هیچ زیرمدول
ناصفری مکمل ندارد (چرا؟).» سپس فقط آغازِ $\mathbb{Q}=M\oplus N$ و سه‌نقطه آمده است.
\evidence{L05-Q05-complement.png}
\textbf{خوانش‌های ممکن:} همهٔ زیرمدول‌های ناصفر، به‌صورت لفظی؛ یا همهٔ
زیرمدول‌های ناصفرِ \textbf{سره}، به‌عنوان اصلاح پیشنهادی.

\textbf{دلیل پرسش:} خود $\mathbb{Q}$ زیرمدولی ناصفر است و مکمل $(0)$ دارد؛
پس صورت لفظی این استثنا را هم در بر می‌گیرد. ادامهٔ اثبات در این فایل نیست و
از بیرون منبع تکمیل نشده است.

\textbf{پرسش:} آیا در اسلاید بنویسم: «در $\mathbb{Q}$ به عنوان $\mathbb{Z}$-مدول،
هیچ زیرمدول ناصفرِ سره‌ای مکمل ندارد»؟
\answer
\newpage

\qtitle{L05-Q06}{فرض حلقه و جهت ایده‌آل در مثال خارج‌قسمت}
\pending
\sourcepage{۳، مثالِ بالای صفحه، پس از ادامهٔ مثال ایده‌آل مینیمال از صفحهٔ ۲}
\textbf{عبارت منبع:} اگر $I$ یک «ایده‌آل ماکسیمال» در حلقهٔ $R$ باشد، آنگاه
$R/I$ به عنوان $R$-مدول ساده است. در این مثال، جابجایی حلقه تکرار نشده و
واژهٔ «چپ» هم در کنار ایده‌آل نیامده است.
\evidence{L05-Q06-maximal-ideal.png}
\textbf{خوانش‌های ممکن:}
\begin{enumerate}
    \item $R$ همچنان حلقه‌ای جابجایی و یک‌دار مانند مثال پیشین است و $I$ ایده‌آل ماکسیمال آن است.
    \item $R$ حلقه‌ای یک‌دار و دلخواه است و منظور از $I$ یک ایده‌آل \textbf{چپ}
    ماکسیمال است، سازگار با قرارداد مدول‌های چپ.
\end{enumerate}
\textbf{دلیل پرسش:} در حلقه‌ای دلخواه، ماکسیمال‌بودنِ دوسویهٔ $I$ به‌تنهایی
ساده‌بودنِ مدول چپ $R/I$ را نتیجه نمی‌دهد.

\textbf{پرسش:} کدام صورت در اسلاید بیاید: حلقهٔ جابجایی و یک‌دار با ایده‌آل
ماکسیمال، یا حلقهٔ یک‌دار دلخواه با ایده‌آل چپ ماکسیمال؟
\answer
\end{document}
